\documentclass[11pt,a4paper,UTF8,fontset=fandol]{ctexart}
\usepackage[margin=24mm]{geometry}
\usepackage{amsmath,amssymb,amsthm,booktabs}
\usepackage[colorlinks=true,linkcolor=black,citecolor=blue,urlcolor=blue]{hyperref}
\hypersetup{pdftitle={单一弱曲率谱值下的梯度重启时间},
 pdfauthor={匿名},pdfsubject={AI 辅助研究工作稿，待作者逐项人工核阅}}
\newtheorem{theorem}{定理}[section]
\newtheorem{lemma}[theorem]{引理}
\newtheorem{proposition}[theorem]{命题}
\DeclareMathOperator{\diag}{diag}
\newcommand{\eps}{\varepsilon}
\newcommand{\norm}[1]{\lVert #1\rVert}
\allowdisplaybreaks[1]
\title{单一弱曲率谱值下的梯度重启时间}
\author{\small AI 辅助研究工作稿\quad 待作者逐项人工核阅}
\date{2026 年 11 月 26 日 \quad v0.1}
\begin{document}
\maketitle
\begin{abstract}
考察正定二次目标上的离散 Nesterov 梯度重启。每段从零动量出发，
取动量系数 \(\beta_k=(k-1)/(k+2)\)。触发判据为正时，舍弃试探点并从
上一个已接受点重算。设步进矩阵的一个谱值 \(s=1-\eps\) 趋于 \(1\)，
其余不同谱值的个数有统一上界，且均留在固定区间
\([\sigma,\rho]\subset(0,1)\)。记 \(n\) 为纯弱曲率残差首次非正的试探
编号，\(K\) 为所有非零段初值中首次触发或到达极小点所需的最多试探
次数。对充分小的 \(\eps\)，一致有 \(K\in\{n,n+1\}\)，且
\(\sqrt{\eps}K\to j_{1,1}\)，其中 \(j_{1,1}\) 是 \(J_1\) 的第一个正零点。
若 \(K=n+1\)，则 \(0\le s_n-s\le C n^{-6}\rho^n\)，
其中 \(s_n\) 是第 \(n\) 步残差的最大正根。对任一固定快谱值
\(s_f\)，存在 \(a>0\)，使二谱模型沿无穷多个 \(n\)，
在区间 \([s_n-a n^{-6}s_f^n,s_n]\) 上达到 \(K=n+1\)。两个有理算例分别给出
\((n,K)=(7,8)\) 和 \((36,37)\)。
\end{abstract}

\section{问题与已有结果}
既有工作稿~\cite{archive} 对固定有限谱和远离端点的紧谱族给出统一的有限返回。本文考察单个步进谱值趋于 \(1\) 时，最坏试探次数如何增长。研究变量是一个逼近 \(1\) 的步进谱值和留在内部区间的快谱值。后者允许随参数变化和互相碰撞。初始点可以属于任意谱面。

O'Donoghue 和 Candès~\cite{oc} 提出了梯度重启判据。本文采用 Bao、Chen、Li 和 Shen~\cite{bao25} 的拒绝试探并从旧点重算的版本，取线性的动量参数序列。他们的命题 2 和推论 2 以统一重启间隔为假设来改进收敛界。Moursi、Pavlovic 和 Vavasis~\cite[命题 3.2]{moursi} 已刻画一维判据的过零机制，他们的重启保留试探点。Kindermann~\cite[定理 1]{kindermann} 给出离散残差的 Gegenbauer 表示，Su、Boyd 和 Candès~\cite{su} 给出二次目标连续极限的 Bessel 表示。本文在这些工具之上控制离散相邻差分和混合谱的严格符号，将最坏时刻限制在两个相邻整数，并确定额外一步的参数尺度。Bao、Chen、Li 和 Sun~\cite[定理 1.1]{bao26} 的永不重启构造属于更广的光滑强凸类，其初始段使用非二次项安排后续相位。

\section{算法与主定理}
设 \(f(x)=x^\top Hx/2\)，其中 \(H=H^\top\succ0\) 为有限维矩阵。取 \(0<h<1/\lambda_{\max}(H)\)，令 \(A=I-hH\)。每段从 \(x_0=y_1=u\ne0\) 出发。第 \(k\) 次试探计算 \(z_k=Ay_k\)。若 \(z_k=0\)，算法终止。否则计算
\[
 D_k=\langle z_k-x_{k-1},\, y_k-z_k\rangle .
\]
当 \(D_k>0\) 时，舍弃 \(z_k\)，从 \(x_{k-1}\) 重算一步 \(Ax_{k-1}\) 后重启。若 \(D_k\le0\)，接受 \(x_k=z_k\)，设置
\[
 \beta_k=\frac{k-1}{k+2},\qquad
 y_{k+1}=x_k+\beta_k(x_k-x_{k-1}).
\]
若在第 \(k\) 次试探触发，该段进行 \(k\) 次试探和一次从旧点的重算，共使用 \(k+1\) 次梯度计算。若第 \(k\) 次试探到达零点，该段使用 \(k\) 次。以下以试探次数计时，并记 \(\tau\) 为首次触发或终止的试探编号。

先形式地沿接受分支计算所有试探，得到 \(z_k=p_k(A)u\)，其中
\begin{equation}
 p_0(v)=1,\quad p_1(v)=v,\quad
 p_{k+1}(v)=v\bigl[p_k(v)+\beta_k(p_k(v)-p_{k-1}(v))\bigr].
 \label{eq:recurrence}
\end{equation}
这里用 \(v\in(0,1)\) 表示任一谱值。令 \(\Delta_k(v)=p_k(v)-p_{k-1}(v)\)。
记 \(P_i\) 为 \(A\) 的谱值 \(s_i\) 对应的整个特征子空间的正交投影，并置 \(w_i=\norm{P_i u}^2/\norm{u}^2\)，则
\begin{equation}
 c_k(v)=\frac{1-v}{v}p_k(v)\Delta_k(v),
 \qquad
 d_k(w)=\frac{D_k}{\norm{u}^2}=\sum_i w_i c_k(s_i).
 \label{eq:criterion}
\end{equation}
这是因为 \(y_k-z_k=(A^{-1}-I)z_k\)，且不同谱方向正交。截至并包括首次触发或终止的那次试探，形式递推给出的试探点与实际计算一致。

\begin{theorem}[最坏时刻至多延迟一步]\label{thm:main}
固定整数 \(m\ge1\) 及区间 \(I=[\sigma,\rho]\)，其中 \(0<\sigma\le\rho<1\)。设 \(A\) 有一个谱值 \(s=1-\eps>\rho\)，其余至多 \(m\) 个不同谱值 \(\boldsymbol{s}\) 位于 \(I\)。对所有充分小的 \(\eps>0\)，定义
\[
 n=n(s)=\min\{k\ge1:p_k(s)\le0\},
 \qquad
 K=K(s,\boldsymbol{s})=\max_{u\ne0}\tau.
\]
有 \(n\le K\le n+1\)，并且
\begin{equation}
 n=\frac{j_{1,1}}{\sqrt{\eps}}+O(1),\qquad
 \lim_{\eps\downarrow0}\sqrt{\eps}\,K=j_{1,1}
       =3.8317059702\ldots .
 \label{eq:limit}
\end{equation}
小参数阈值和极限关于 \(I\) 内的快谱值一致。该结论包含快谱值固定、随 \(\eps\) 变化、互相重合及对应初始权重趋零的情形。
\end{theorem}

在这个谱族中，\(\eps=h\lambda_{\min}(H)\)。其余不同特征值满足 \(h\lambda_i\in[1-\rho,1-\sigma]\)。如取 \(h\) 固定，这描述 \(\lambda_{\min}(H)\) 单独趋零。若固定有多个不同谱值的 \(H\) 并令 \(h\downarrow0\)，则所有步进谱值都趋近 \(1\)，属于另一种渐近问题。同一个 \(H,h\) 的每次重启都从零动量开始，所以本定理也给这些段的统一试探次数上界。按上面的重算口径，最坏梯度调用数介于 \(K\) 与 \(K+1\) 之间，也满足式 \eqref{eq:limit} 的首项极限。

\section{弱曲率模态的残差与过零}
\begin{lemma}[端点附近的相邻差分]\label{lem:bessel}
写 \(s=1-t^2\)、\(t>0\)，并定义
\[
 F(x)=\frac{2J_1(x)}{x}
   =\sum_{\ell=0}^\infty\frac{(-1)^\ell(x^2/4)^\ell}{\ell!(\ell+1)!},
 \qquad F(0)=1 .
\]
若 \(k\ge1\) 且 \((k+\tfrac12)t\le M\)，其中 \(M>0\) 固定，则当 \(t\downarrow0\) 时有一致估计
\begin{align}
 s^{-(k+1)/2}p_k(s)&=F((k+\tfrac12)t)+O_M(t^2),\label{eq:qexp}\\
 \Delta_k(s)&=tF'(kt)+O_M(t^2)\quad (k\ge2).\label{eq:diffexp}
\end{align}
\end{lemma}
\begin{proof}
残差公式~\cite[定理 1]{kindermann} 在参数 \(2\) 时，经 Gegenbauer 和 Jacobi 对应关系写成
\begin{equation}
 p_k(v)=v^{(k+1)/2}\frac{P_{k-1}^{(1,1)}(\sqrt{v})}{k}.
 \label{eq:jacobi}
\end{equation}
也可用 Jacobi 三项递推直接验证式 \eqref{eq:recurrence}。
设
\[
 \nu_k=k+\tfrac12,\qquad
 \zeta_k(v)=\nu_k\sqrt{2(1-\sqrt{v})},\qquad
 b_{k,\ell}=
 \begin{cases}
 \prod_{r=0}^{\ell-1}\left(1-\dfrac{(r+3/2)^2}{\nu_k^2}\right),&0\le\ell<k,\\[3pt]
 0,&\ell\ge k .
 \end{cases}
\]
其中空乘积为 \(1\)。由 Jacobi 的超几何式~\cite[式 18.5.7]{dlmf}，
\begin{equation}
 \frac{P_{k-1}^{(1,1)}(\sqrt v)}{k}
   =Q_k(\zeta_k(v)),\qquad
 Q_k(x)=\sum_{\ell\ge0}
   \frac{(-1)^\ell b_{k,\ell}(x^2/4)^\ell}{\ell!(\ell+1)!}.
 \label{eq:product}
\end{equation}
系数乘积来自恒等式
\((k-1-r)(k+2+r)=\nu_k^2-(r+3/2)^2\)。当 \(\ell<k\) 时，乘积各项位于 \([0,1]\)。当 \(\ell\ge k\) 时，上式的第一项零因子对应 \(r=k-1\)。因此对 \(\ell\ge1\)，
\[
 0\le1-b_{k,\ell}
 \le \nu_k^{-2}\sum_{r=0}^{\ell-1}(r+3/2)^2 .
\]
对任意固定 \(R\)，用这个估计比较两组级数及其逐项导数。在 \(|x|\le R\) 上，带有额外三次或四次 \(\ell\) 因子的级数仍绝对一致收敛，给出
\begin{equation}
 |Q_k(x)-F(x)|\le C_R x^2/\nu_k^2,\qquad
 |Q_k'(x)-F'(x)|\le C_R/\nu_k^2.
 \label{eq:seriesbounds}
\end{equation}
第一个上界保留了 \(\ell\ge1\) 提供的 \(x^2\) 因子。对 \(v=s=1-t^2\)，有
\(\zeta_k(s)=\nu_k t \sqrt{2/(1+\sqrt{s})}=\nu_k t(1+O(t^2))\)。
在 \(\nu_k t\le M\) 时，式 \eqref{eq:seriesbounds} 和 \(F\) 的有界导数给出式 \eqref{eq:qexp}。再令
\(q_j=s^{-(j+1)/2}p_j(s)\)。对 \(k\ge2\)，
\[
 \Delta_k(s)=s^{k/2}\big[\sqrt{s}\,q_k-q_{k-1}\big].
\]
式 \eqref{eq:qexp} 在相邻两个指标上合用，给出
\(q_k-q_{k-1}=tF'(kt)+O_M(t^2)\)。
因为 \(\sqrt{s}-1=O(t^2)\)、\(q_k=O_M(1)\)，且 \(s^{k/2}=1+O_M(t)\)，得到式 \eqref{eq:diffexp}。
\end{proof}

\begin{lemma}[单模态首次试探事件]\label{lem:cross}
令 \(j_*=j_{1,1}\)。当 \(t=\sqrt{1-s}\downarrow0\) 时，\(n(s)\) 存在，并满足
\begin{equation}
 |nt-j_*|\le \tfrac12t+O(t^2).
 \label{eq:round}
\end{equation}
在纯 \(s\) 模态上首次试探事件恰为 \(n\)。另外存在 \(c>0\)，使所有充分小的 \(t\) 都有
\begin{equation}
 p_{n+1}(s)\le-ct,\qquad
 \Delta_{n+1}(s)\le-ct,\qquad
 c_{n+1}(s)\ge c^2t^4 .
 \label{eq:late}
\end{equation}
\end{lemma}
\begin{proof}
\(F\) 在 \([0,j_*)\) 上严格为正，在 \(j_*\) 有简单零点，故 \(F'(j_*)<0\)~\cite[第 10.21 节]{dlmf}。取足够小的 \(\delta>0\)，使 \(F\) 在 \((j_*,j_*+\delta]\) 上严格为负。式 \eqref{eq:qexp} 对 \(k t\) 落于固定紧区间内一致成立，它排除在 \((k+1/2)t\le j_*-\delta\) 时出现非正值。取 \((k+1/2)t\) 逼近 \(j_*+\delta\) 又保证出现负值。由此 \(nt\to j_*\)。相邻两个指标 \(n,n-1\) 的正规化残差分别非正和严格为正，且 \(F'\) 在 \(j_*\) 附近严格为负且离开零。再用式 \eqref{eq:qexp} 得到
\((n+1/2)t\ge j_*-O(t^2)\) 与 \((n-1/2)t\le j_*+O(t^2)\)，即式 \eqref{eq:round}。

在 \(k\le n\) 且 \(p_{k-1}>0\) 时，式 \eqref{eq:recurrence} 给出
\(\Delta_1=-t^2\) 和
\(\Delta_k=-t^2p_{k-1}+s\beta_{k-1}\Delta_{k-1}<0\)（\(k\ge2\)）。
于是 \(k<n\) 的判据均严格负。在 \(k=n\)，若 \(p_n<0\)，判据严格正；若 \(p_n=0\)，试探到达零点。最后，式 \eqref{eq:diffexp} 在 \(k=n+1\) 给出 \(\Delta_{n+1}\le-ct\)。将其加到 \(p_n\le0\) 上，再代入式 \eqref{eq:criterion}，即得式 \eqref{eq:late}。
\end{proof}

\section{对混合谱的统一控制}
本节先用已有的快谱正裕量方法~\cite[第 3 节]{archive}限制快谱权重。接着利用晚期的指数衰减，证明抵消只能使最坏首次事件晚一步。
\begin{lemma}[快谱的早期正裕量与晚期衰减]\label{lem:fast}
固定 \(m\) 和 \(I=[\sigma,\rho]\subset(0,1)\)。存在 \(N_0\)、\(\eta>0\) 与 \(C>0\)，使对任意至多 \(m\) 个位于 \(I\) 的谱值 \(v_i\) 及任意归一化权重 \(\widehat w\)，有
\begin{equation}
 \max_{1\le k\le N_0}\sum_i\widehat w_i c_k(v_i)\ge\eta,\qquad
 |c_k(v)|\le C\rho^k k^{-3}\quad(v\in I,\ k\ge1).
 \label{eq:fast}
\end{equation}
另外，对任一固定 \(v\in(0,1)\)，有常数 \(b_v>0\) 和无穷多个 \(k\)，使
\begin{equation}
 c_k(v)\le-b_v v^k k^{-3}.
 \label{eq:fastnegative}
\end{equation}
\end{lemma}
\begin{proof}
令 \(\theta=\arccos\sqrt v\)，\(\phi_k=(k+\tfrac12)\theta-3\pi/4\)。在 \(I\) 上，Jacobi Darboux 展开~\cite[式 18.15.4\_5]{dlmf} 与式 \eqref{eq:jacobi} 给出
\[
 p_k(v)=a_v v^{(k+1)/2} k^{-3/2}\bigl(\cos\phi_k+O_I(k^{-1})\bigr),
 \qquad a_v=\sqrt{8/\pi}\,(\sin\theta)^{-3/2}.
\]
提出 \(p_k-p_{k-1}\) 的共同因子后，主项使用
\(\sqrt v\cos\phi_k-\cos(\phi_k-\theta)=-\sin\theta\sin\phi_k\)。
将两式相乘得
\[
 c_k(v)=-\frac{(1-v)a_v^2\tan\theta}{2}
 v^k k^{-3}\bigl(\sin(2\phi_k)+O_I(k^{-1})\bigr).
\]
这给出式 \eqref{eq:fast} 的晚期估计。对每个固定 \(v\)，若 \(\theta/\pi\) 无理，相位模 \(2\pi\) 稠密；若有理，其既约分母至少为 \(3\)，周期内 \(2\phi_k\bmod 2\pi\) 为等间距点，正弦和为零，且至多两个点的正弦为零。两个情形均有严格正和严格负、并分别远离零的正弦子列。因此得到式 \eqref{eq:fastnegative}，也得到相应的正判据子列。

任给快谱参数和权重，合并相等的谱值。沿最大活跃谱值的正判据子列，该模态的正项以严格指数比支配其余各项，所以某个加权判据严格为正。补零权至 \(m\) 个位置，参数空间 \(I^m\) 与归一化权重单纯形的乘积是紧集。连续函数
\(\max_{1\le k\le N}\sum_i\widehat w_i c_k(v_i)\) 的严格正集合随 \(N\) 增大，并覆盖紧集，故某个 \(N_0\) 已覆盖。该函数在紧集上的最小值严格为正，给出 \(\eta\)。
\end{proof}

\begin{proof}[定理 \ref{thm:main} 的证明]
令 \(q\) 为快谱总权重，\(1-q\) 为慢谱权重。当 \(q>0\)，快谱的归一化权重记为 \(\widehat w\)。因 \(p_k(1)=1\)，固定前 \(N_0\) 个指标都有 \(|c_k(s)|\le C_0t^4\)，其中 \(t=\sqrt\eps\)。若实际轨道在前 \(N_0\) 次尚未发生试探事件，则由引理 \ref{lem:fast} 选一个 \(k_0\le N_0\) 使快谱判据至少为 \(\eta\)。在这个试探有
\[
 0\ge d_{k_0}\ge -(1-q)C_0t^4+q\eta .
\]
于是 \(q<1\)，且
\begin{equation}
 \omega:=q/(1-q)\le (C_0/\eta)t^4.
 \label{eq:weight}
\end{equation}
当 \(q=0\) 时，式 \eqref{eq:weight} 自动成立。取 \(t\) 足够小，使 \(n>N_0\)。
若到第 \(n+1\) 次试探仍无更早事件，则
\[
 \frac{d_{n+1}}{1-q}
 \ge c^2t^4-\omega C\rho^{n+1}(n+1)^{-3}>0
\]
由式 \eqref{eq:late}、\eqref{eq:fast} 和 \eqref{eq:weight} 得到。
同时慢谱活跃且 \(p_{n+1}(s)\ne0\)，故这一步确为触发。
纯慢谱方向在第 \(n\) 次发生事件，所以至少有一个初值达到 \(n\)。
所有 \(\tau\) 是不超过 \(n+1\) 的正整数，故最大值存在。引理 \ref{lem:cross} 给出其一致极限。
\end{proof}

\section{额外一步的参数区间}
对 \(n\ge3\)，令 \(s_n\) 为 \(P_{n-1}^{(1,1)}\) 最大正根的平方，并置 \(s_2=0\)。Jacobi 多项式的根在 \((-1,1)\) 内简单且严格交错~\cite[第 18.2(vi) 节]{dlmf}。由式 \eqref{eq:jacobi}，当 \(s\in(s_{n-1},s_n]\) 时，\(p_1,\ldots,p_{n-1}\) 均严格为正，\(p_n\le0\)，从而 \(n(s)=n\)。这些区间覆盖 \((0,1)\)。事实上，式 \eqref{eq:qexp} 在 \(t=(j_*+\delta)/(n+1/2)\) 给出大 \(n\) 时的一个负值，故交错的最大根 \(s_n\) 趋于 \(1\)。在 \(s=s_n\) 时，式 \eqref{eq:qexp} 和 \(F'(j_*)<0\) 进一步给出
\begin{equation}
 1-s_n=\frac{j_*^2}{(n+1/2)^2}+O(n^{-4}),
 \qquad
 s_n-s_{n-1}=2j_*^2n^{-3}+O(n^{-4}).
 \label{eq:rootshift}
\end{equation}

\begin{proposition}[延迟区间的上下界]\label{prop:window}
在定理 \ref{thm:main} 的谱族中，若 \(K=n+1\)，则
\begin{equation}
 0\le s_n-s\le C n^{-6}\rho^n
 \label{eq:upperwindow}
\end{equation}
对充分大的 \(n=n(s)\) 成立，常数可在同一个快谱族中统一。反之，固定任意一个快谱值 \(s_f\in(0,1)\)。存在 \(a>0\) 和无穷多个 \(n\)，使对所有
\begin{equation}
 s\in[s_n-a n^{-6}s_f^n,s_n]
 \label{eq:lowerwindow}
\end{equation}
的二谱模型都有 \(n(s)=n\) 且 \(K=n+1\)。
\end{proposition}
\begin{proof}
先给出统一的标量估计。在整个区间 \(s\in(s_{n-1},s_n]\) 上，\(\sqrt{1-s}\asymp n^{-1}\)、\(\zeta_n(s)\to j_*\)，其中 \(\asymp\) 表示上下各有一个严格正常数。对式 \eqref{eq:product} 求导并使用式 \eqref{eq:seriesbounds}，有
\[
 Q_n'(\zeta_n(s))\to F'(j_*)<0,\qquad
 -\zeta_n'(s)=\frac{n+1/2}{2\sqrt{s}\sqrt{2(1-\sqrt{s})}}
 \asymp n^2 .
\]
正规化残差在 \(s_n\) 恰为零，所以中值定理给
\(\lvert p_n(s)\rvert\asymp n^2(s_n-s)\)，其中
\(s^{(n+1)/2}\to1\) 保留同阶性。
在 \(s=s_n\) 时两边同为零。式 \eqref{eq:diffexp} 给
\(\Delta_n(s)<0\) 且 \(\lvert\Delta_n(s)\rvert\asymp n^{-1}\)。
综合得到
\begin{equation}
 c_n(s)\asymp n^{-1}(s_n-s).
 \label{eq:cn-scale}
\end{equation}
若某个初值实现 \(K=n+1\)，则截至 \(n\) 无事件，并由式 \eqref{eq:weight} 有 \(\omega\le Cn^{-4}\)。在第 \(n\) 次，以慢谱权重归一化，非正判据及式 \eqref{eq:fast} 给出
\(c_n(s)\le \omega Cn^{-3}\rho^n\le C'n^{-7}\rho^n\)。
式 \eqref{eq:cn-scale} 给出式 \eqref{eq:upperwindow}。

证明下界时，由式 \eqref{eq:fastnegative} 取无穷子列，使
\(c_n(s_f)\le-b n^{-3}s_f^n\)，\(b>0\) 固定。由式 \eqref{eq:rootshift}，对该子列中充分大的 \(n\)，式 \eqref{eq:lowerwindow} 的区间落在 \((s_{n-1},s_n]\)。在其中有
\(p_n(s)=O(n^{-4}s_f^n)\) 且
\(p_{n-1}(s)=p_n(s)-\Delta_n(s)\ge c_1/n\)，对任意固定 \(a\) 一致。对 \(1\le k\le n/2\)，式 \eqref{eq:qexp} 给 \(p_k(s)\ge c_2>0\)，因为这些自变量留在 \(F\) 的第一零点之前，并与其保持正距离。对 \(n/2<k<n\)，由首次过零前的单调性知 \(p_k(s)\ge p_{n-1}(s)\ge c_1/n\)。又由递推有 \(-\Delta_k(s)\ge (1-s)p_{k-1}(s)\) 对全部 \(1\le k<n\) 成立。于是
\[
 -c_k(s)\ge \frac{(1-s)^2}{s}p_{k-1}(s)p_k(s)
 \ge
 \begin{cases}
  c_3n^{-4},&1\le k\le n/2,\\
  c_4n^{-6},&n/2<k<n.
 \end{cases}
\]
对每个使 \(c_k(s_f)>0\) 的 \(k<n\)，以上两式与
\(\lvert c_k(s_f)\rvert\le C s_f^k k^{-3}\) 表明
\(-c_k(s)/c_k(s_f)\ge c_5 n^{-4}\)，其中晚期使用 \(s_f^{n/2}\) 的指数衰减。
因此，令快谱权与慢谱权之比 \(\omega=\gamma n^{-4}\)，其中 \(0<\gamma<c_5\)，可保持所有 \(k<n\) 的判据严格负。常数 \(c_1,\ldots,c_5\) 可先在 \(a\le1\) 下统一选取。取 \(a\) 进一步满足 \(Ca<\gamma b\)，由式 \eqref{eq:cn-scale} 的上界，
\[
 c_n(s)+\omega c_n(s_f)\le (Ca-\gamma b)n^{-7}s_f^n<0 .
\]
第 \(n\) 次试探的快分量仍非零，因为 \(c_n(s_f)<0\) 且快谱权为正。这也覆盖 \(s=s_n\) 的恰好过零情形。前 \(n-1\) 次慢分量始终非零。故确实在第 \(n+1\) 次才第一次触发，完成下界。
\end{proof}

\section{二谱的判定与有理证书}
设只有一个快谱值 \(s_f\)，且 \(s\) 足够近 \(1\)、\(p_n(s)<0\)。
写 \(A_k=c_k(s)\)、\(B_k=c_k(s_f)\)。前 \(n-1\) 次有 \(A_k<0\)，第 \(n\) 次有 \(A_n>0\)。对快谱权、慢谱权之比 \(\omega\) 而言，前 \(n-1\) 次未触发恰等价于 \(\omega\le U\)，其中
\[
 U=\min_{\substack{1\le k<n\\ B_k>0}}\frac{-A_k}{B_k}.
\]
小参数下，由快谱的早期正裕量知集合非空，纯快谱在第 \(n\) 次试探之前已有首次事件。第 \(n\) 次仍未触发恰等价于 \(B_n<0\) 且 \(\omega\ge A_n/(-B_n)\)。
此时 \(p_n(s)<0\) 也排除提前终止。因此
\[
 K=n+1\quad\Longleftrightarrow\quad
 B_n<0\ \text{且}\ A_n/(-B_n)\le U .
\]
等号可实现，因为判据为零时继续接受试探。若 \(p_n(s)=0\)，则纯慢方向在第 \(n\) 次终止，对混合方向仍可取 \(0<\omega<U\) 使第 \(n\) 次之前的判据严格负。在第 \(n\) 次，混合试探点为零的充要条件是 \(p_n(s_f)=0\)，其判据与 \(B_n\) 同号。因此在 \(p_n(s)=0\) 的情形，\(K=n+1\) 的充要条件是 \(p_n(s_f)\ne0\) 且 \(B_n\le0\)。这包含命题 \ref{prop:window} 的区间右端点。

下表取 \(h=1\)、\(H=\diag(1-s,1/2)\)。两例固定快谱值 \(s_f=1/2\)，各行的初始点给出达到所报上界的真实轨道。
\begin{center}
\setlength{\tabcolsep}{4pt}
\begin{tabular}{cccccc}
\toprule
 \(s\) & 达到 \(K\) 的 \(u\) & 单模态 \(n\) & 最坏次数 \(K\) & 证书试探 & 系数 \(M\) \\
\midrule
 \(7599/10000\) & \((1,1)^\top\) & \(7\) & \(8\) & \(5,8\) & \(32\) \\
 \(49450896852043/50000000000000\) & \((2,1)^\top\) & \(36\) & \(37\) & \(5,37\) & \(64\) \\
\bottomrule
\end{tabular}
\end{center}
上界证书为 \(c_5(v)+M c_K(v)>0\)，对表中各行的两个谱值 \(v\) 分别成立。若第 \(K\) 次试探后仍无事件，第 5 次和第 \(K\) 次判据均非正。将逐谱值的严格式按初始谱权求和，得到矛盾，故上界覆盖全部初值。下界以指定初值逐步验证第 \(1\) 至 \(K-1\) 次判据严格负，第 \(K\) 次严格正，且各试探均非零。对第一行再取 \(\omega=A_7/(-B_7)\)，可令第 \(7\) 次判据恰为零而试探非零，直到第 \(8\) 次触发。在另一标量算例中取 \(s=1/5\)，第三次试探精确到达零点。精确复算使用有理数，独立地按标量残差递推及直接坐标算法核对，并保留判据和证书系数。渐近结论由前文证明给出。

本文处理固定数量谱节点、单一弱曲率模态与每个零动量段的首次事件。多个谱值共同逼近 \(1\)、连续谱和按实际段间权重演化计算的尖锐长期收敛率，仍属待证问题。

\paragraph{生成式 AI 使用说明。}
生成式 AI 辅助了来源核对、推导、撰写和算例检查。稿件待作者逐项人工核阅。

\begin{thebibliography}{9}
\bibitem{archive}
《离散梯度重启的有限返回与投影 BB1 的临界条件数收敛》，
AI 辅助公开工作稿 v0.3.0，2026 年 11 月 25 日。
\href{https://taitai2025.github.io/assets/documents/research/restart-projected-bb.pdf}{公开版本}。

\bibitem{oc}
B. O'Donoghue and E. J. Candès,
\emph{Adaptive Restart for Accelerated Gradient Schemes},
Found. Comput. Math. 15 (2015), 715--732.
\href{https://arxiv.org/abs/1204.3982}{arXiv:1204.3982}.

\bibitem{bao25}
C. Bao, L. Chen, J. Li and Z. Shen,
\emph{Accelerated Gradient Methods with Gradient Restart: Global Linear Convergence},
\href{https://arxiv.org/html/2401.07672v2}{arXiv:2401.07672v2}, 2025.

\bibitem{moursi}
W. M. Moursi, V. Pavlovic and S. A. Vavasis,
\emph{Accelerated Gradient Descent: A Guaranteed Bound for a Heuristic Restart Strategy},
\href{https://arxiv.org/html/2310.07674v1}{arXiv:2310.07674v1}, 2023.

\bibitem{kindermann}
S. Kindermann,
\emph{Optimal-order convergence of Nesterov acceleration for linear ill-posed problems},
Inverse Problems 37 (2021), 065002.
\href{https://arxiv.org/html/2101.08168v1}{arXiv:2101.08168v1}.

\bibitem{su}
W. Su, S. Boyd and E. J. Candès,
\emph{A Differential Equation for Modeling Nesterov's Accelerated Gradient Method: Theory and Insights},
J. Mach. Learn. Res. 17 (2016), 1--43.
\href{https://jmlr.org/papers/v17/15-084.html}{期刊全文}。

\bibitem{bao26}
C. Bao, L. Chen, J. Li and D. Sun,
\emph{On the Nonoccurrence of Gradient Restart in Nesterov Accelerated Gradient Method},
\href{https://arxiv.org/html/2610.26546v1}{arXiv:2610.26546v1}, 2026.

\bibitem{dlmf}
NIST Digital Library of Mathematical Functions，
\href{https://dlmf.nist.gov/10.21}{第 10.21 节}，
\href{https://dlmf.nist.gov/18.2.vi}{第 18.2(vi) 节}，
\href{https://dlmf.nist.gov/18.5.E7}{式 18.5.7}，
\href{https://dlmf.nist.gov/18.11.ii}{第 18.11(ii) 节}，
\href{https://dlmf.nist.gov/18.15.E4_5}{式 18.15.4\_5}。
\end{thebibliography}

\end{document}
